\paragraph{escort 温度流的总变差压缩、散度闭式与 Fisher 平直化}
在末位充分化已经把 escort 温度族压缩到一维 Bernoulli 参数之后，总变差塌缩、Rényi--KL 散度几何与 Fisher 度量都可沿同一末位两块分解完全显式化。

\begin{theorem}[escort 温度族在总变差意义下塌缩为一维 Bernoulli 指数族]\label{thm:xi-foldbin-escort-tv-collapse-onebit-gibbs}
对任意 \(q\ge 0\)，定义 escort 分布
\[
\pi_{m,q}(w):=\frac{(d_m^{\mathrm{bin}}(w))^q}{S_q^{\mathrm{bin}}(m)},
\qquad
S_q^{\mathrm{bin}}(m):=\sum_{u\in X_m}(d_m^{\mathrm{bin}}(u))^q.
\]
并按末位将 \(X_m\) 分成
\[
A_0:=\{w\in X_m:\ w_m=0\},
\qquad
A_1:=\{w\in X_m:\ w_m=1\}.
\]
记
\[
\theta(q):=\frac{1}{1+\varphi^{q+1}},
\]
并定义块内均匀近似
\[
\bar\pi_{m,q}(w):=
\frac{1-\theta(q)}{|A_0|}\mathbf 1_{A_0}(w)
+
\frac{\theta(q)}{|A_1|}\mathbf 1_{A_1}(w).
\]
则对每个固定 \(q\ge 0\)，存在常数 \(C_q>0\) 使
\[
\|\pi_{m,q}-\bar\pi_{m,q}\|_{\mathrm{TV}}
\le
C_q\Bigl(\frac{\varphi}{2}\Bigr)^m.
\]
进一步，\(\bar\pi_{m,q}\) 可写成稳定层均匀分布 \(u_m(w):=1/|X_m|\) 的单参数指数倾斜
\[
\bar\pi_{m,q}(w)=Z_{m,q}^{-1}u_m(w)e^{-\beta_{m,q}w_m},
\]
且
\[
\beta_{m,q}\longrightarrow q\log\varphi.
\]
因此整个 escort 温度族在全变差精度下严格塌缩为由末位比特 \(w_m\) 支配的一维 Bernoulli 指数族。
\end{theorem}
\begin{proof}
由定理 \ref{thm:fold-bin-two-state-asymptotic}，对 \(i\in\{0,1\}\) 与 \(w\in A_i\) 一致有
\[
d_m^{\mathrm{bin}}(w)
=
\left(\frac{2}{\varphi}\right)^m
\varphi^{-i}
\Bigl(1+O\!\Bigl(\Bigl(\frac{\varphi}{2}\Bigr)^m\Bigr)\Bigr).
\]
对固定 \(q\ge 0\) 取 \(q\) 次幂可得
\[
(d_m^{\mathrm{bin}}(w))^q
=
\left(\frac{2}{\varphi}\right)^{qm}
\varphi^{-qi}
\Bigl(1+O_q\!\Bigl(\Bigl(\frac{\varphi}{2}\Bigr)^m\Bigr)\Bigr),
\]
且误差在块内一致。于是
\[
\sum_{w\in A_i}(d_m^{\mathrm{bin}}(w))^q
=
|A_i|
\left(\frac{2}{\varphi}\right)^{qm}
\varphi^{-qi}
\Bigl(1+O_q\!\Bigl(\Bigl(\frac{\varphi}{2}\Bigr)^m\Bigr)\Bigr).
\]
再用 \(|A_0|=F_{m+1}\)、\(|A_1|=F_m\) 以及
\(
F_m/F_{m+1}=\varphi^{-1}+O(\varphi^{-2m})
\)，得到
\[
\pi_{m,q}(A_1)=\theta(q)+O_q\!\Bigl(\Bigl(\frac{\varphi}{2}\Bigr)^m\Bigr),
\qquad
\pi_{m,q}(A_0)=1-\theta(q)+O_q\!\Bigl(\Bigl(\frac{\varphi}{2}\Bigr)^m\Bigr).
\]
另一方面，对 \(w\in A_i\) 有
\[
\pi_{m,q}(w)
=
\frac{\pi_{m,q}(A_i)}{|A_i|}
\Bigl(1+O_q\!\Bigl(\Bigl(\frac{\varphi}{2}\Bigr)^m\Bigr)\Bigr).
\]
将其与 \(\bar\pi_{m,q}(w)\) 比较，并在两块上分别求和，即得
\[
\sum_{w\in X_m}\bigl|\pi_{m,q}(w)-\bar\pi_{m,q}(w)\bigr|
=
O_q\!\Bigl(\Bigl(\frac{\varphi}{2}\Bigr)^m\Bigr),
\]
从而给出总变差界。

对指数倾斜表示，只需令
\[
e^{-\beta_{m,q}}
:=
\frac{\bar\pi_{m,q}(A_1)}{\bar\pi_{m,q}(A_0)}
\cdot
\frac{|A_0|}{|A_1|}
=
\frac{\theta(q)}{1-\theta(q)}
\cdot
\frac{F_{m+1}}{F_m}.
\]
由于
\(
\theta(q)/(1-\theta(q))=\varphi^{-(q+1)}
\)
且 \(F_{m+1}/F_m\to\varphi\)，故
\[
e^{-\beta_{m,q}}\to \varphi^{-q},
\]
即 \(\beta_{m,q}\to q\log\varphi\)。证毕。
\end{proof}
\leanverified{paper_xi_foldbin_escort_tv_collapse_onebit_gibbs}

\begin{theorem}[escort 参数在总变差上保持严格可辨识]\label{thm:xi-foldbin-escort-tv-identifiable-limit-distance}\leanverified{paper\_xi\_foldbin\_escort\_tv\_identifiable\_limit\_distance}
沿用定理 \ref{thm:xi-foldbin-escort-tv-collapse-onebit-gibbs} 的记号。则对任意固定 \(p,q\ge 0\)，有
\[
\|\pi_{m,q}-\pi_{m,p}\|_{\mathrm{TV}}
=
|\theta(q)-\theta(p)|
+O_{p,q}\!\Bigl(\Bigl(\frac{\varphi}{2}\Bigr)^m\Bigr),
\]
从而
\[
\lim_{m\to\infty}\|\pi_{m,q}-\pi_{m,p}\|_{\mathrm{TV}}
=
|\theta(q)-\theta(p)|.
\]
特别地，极限参数可由单一标量 \(\theta(q)\) 反演为
\[
q=\log_\varphi\!\left(\frac{1-\theta(q)}{\theta(q)}\right)-1.
\]
\end{theorem}
\begin{proof}
由三角不等式，
\[
\Bigl|
\|\pi_{m,q}-\pi_{m,p}\|_{\mathrm{TV}}
-\|\bar\pi_{m,q}-\bar\pi_{m,p}\|_{\mathrm{TV}}
\Bigr|
\le
\|\pi_{m,q}-\bar\pi_{m,q}\|_{\mathrm{TV}}
+\|\pi_{m,p}-\bar\pi_{m,p}\|_{\mathrm{TV}},
\]
而右端由定理 \ref{thm:xi-foldbin-escort-tv-collapse-onebit-gibbs} 为
\(
O_{p,q}((\varphi/2)^m)
\)。

另一方面，\(\bar\pi_{m,q}\) 与 \(\bar\pi_{m,p}\) 都在同一分块 \(A_0\sqcup A_1\) 上块内均匀，且两块质量分别为
\(
(1-\theta(q),\theta(q))
\)
与
\(
(1-\theta(p),\theta(p))
\)。
因此
\[
\|\bar\pi_{m,q}-\bar\pi_{m,p}\|_{\mathrm{TV}}
=
|\theta(q)-\theta(p)|.
\]
代回即得所示展开与极限。
最后由
\(
\theta(q)=1/(1+\varphi^{q+1})
\)
直接解得反演公式。证毕。
\end{proof}

\begin{theorem}[escort 温度族之间的 Rényi--KL 几何在极限中封闭为显式二点公式]\label{thm:xi-foldbin-escort-renyi-kl-two-point-closure}
沿用定理 \ref{thm:xi-foldbin-escort-tv-collapse-onebit-gibbs} 的记号。对 \(q,r\ge 0\) 与
\(\alpha>0\)、\(\alpha\neq 1\)，定义 Rényi 散度
\[
D_\alpha(\pi_{m,q}\Vert \pi_{m,r})
:=
\frac{1}{\alpha-1}
\log\sum_{w\in X_m}\pi_{m,q}(w)^\alpha \pi_{m,r}(w)^{1-\alpha}.
\]
则极限
\[
\lim_{m\to\infty}D_\alpha(\pi_{m,q}\Vert \pi_{m,r})
=:\mathcal D_\alpha(q,r)
\]
存在，且满足
\[
\mathcal D_\alpha(q,r)
=
\frac{1}{\alpha-1}
\log\!\Bigl(
(1-\theta(q))^\alpha(1-\theta(r))^{1-\alpha}
+
\theta(q)^\alpha\theta(r)^{1-\alpha}
\Bigr).
\]
等价地，
\[
\mathcal D_\alpha(q,r)
=
\frac{1}{\alpha-1}
\log
\frac{\varphi+\varphi^{-\alpha q-(1-\alpha)r}}
{(\varphi+\varphi^{-q})^\alpha(\varphi+\varphi^{-r})^{1-\alpha}}.
\]
进一步，令 \(\alpha\to 1\) 得到极限 KL 散度的闭式
\[
\lim_{m\to\infty}D_{\mathrm{KL}}(\pi_{m,q}\Vert \pi_{m,r})
=
\log\!\frac{1+\varphi^{r+1}}{1+\varphi^{q+1}}
+
\frac{\varphi^{q+1}}{1+\varphi^{q+1}}(q-r)\log\varphi.
\]
因此 escort 温度族的双参数散度几何在极限中完全退化为由 \(\theta(q)\) 决定的二点统计流形。
\end{theorem}
\begin{proof}
记
\(
\varepsilon_m:=(\varphi/2)^m
\)。
由定理 \ref{thm:xi-foldbin-escort-tv-collapse-onebit-gibbs} 的证明，在 \(A_i\) 上有一致估计
\[
\pi_{m,q}(w)
=
\frac{\pi_{m,q}(A_i)}{|A_i|}
\Bigl(1+O_q(\varepsilon_m)\Bigr),
\qquad
\pi_{m,r}(w)
=
\frac{\pi_{m,r}(A_i)}{|A_i|}
\Bigl(1+O_r(\varepsilon_m)\Bigr).
\]
于是
\[
\sum_{w\in A_i}\pi_{m,q}(w)^\alpha\pi_{m,r}(w)^{1-\alpha}
=
\pi_{m,q}(A_i)^\alpha\pi_{m,r}(A_i)^{1-\alpha}
\Bigl(1+O_{q,r,\alpha}(\varepsilon_m)\Bigr),
\]
因为块内的 \(|A_i|\) 因子恰好抵消。

另一方面，由定理 \ref{thm:xi-foldbin-escort-tv-collapse-onebit-gibbs}，
\[
\bigl|\pi_{m,q}(A_1)-\theta(q)\bigr|
\le
\|\pi_{m,q}-\bar\pi_{m,q}\|_{\mathrm{TV}}
=
O_q(\varepsilon_m),
\]
且同理
\(
\pi_{m,q}(A_0)=1-\theta(q)+O_q(\varepsilon_m)
\)；
对 \(r\) 亦然。
因此把 \(i=0,1\) 两块相加，并令 \(m\to\infty\)，即可得到
\[
\mathcal D_\alpha(q,r)
=
\frac{1}{\alpha-1}
\log\!\Bigl(
(1-\theta(q))^\alpha(1-\theta(r))^{1-\alpha}
+
\theta(q)^\alpha\theta(r)^{1-\alpha}
\Bigr).
\]
再代入
\(
\theta(q)=1/(1+\varphi^{q+1})
\)
与
\(
1-\theta(q)=\varphi^{q+1}/(1+\varphi^{q+1})
\)
即可化简为第二种写法。
最后令 \(\alpha\to 1\) 即得 KL 闭式。证毕。
\end{proof}

\begin{theorem}[escort 温度流形的 Fisher 度量是全局平直的]\label{thm:xi-foldbin-escort-logistic-fisher-geometry}
定义 escort 家族的 Fisher 信息
\[
\mathcal I_m(q)
:=
\sum_{w\in X_m}\pi_{m,q}(w)\Bigl(\partial_q\log \pi_{m,q}(w)\Bigr)^2.
\]
则对每个固定 \(q\ge 0\)，有极限
\[
\mathcal I_m(q)\longrightarrow \mathcal I_\infty(q)
:=
(\log\varphi)^2\frac{\varphi^{q+1}}{(1+\varphi^{q+1})^2}.
\]
更强地，极限度量
\[
ds^2=\mathcal I_\infty(q)\,\dd q^2
\]
在变量
\[
u(q):=2\arctan\!\bigl(\varphi^{(q+1)/2}\bigr)
\]
下化为 Euclidean 形式
\[
ds^2=du^2.
\]
因此对任意 \(q_1,q_2\ge 0\)，极限 Fisher 距离都可显式写为
\[
d_F(q_1,q_2)
=
\bigl|u(q_2)-u(q_1)\bigr|
=
2\left|
\arctan\!\bigl(\varphi^{(q_2+1)/2}\bigr)
-
\arctan\!\bigl(\varphi^{(q_1+1)/2}\bigr)
\right|.
\]
特别地，半轴 \(q\in[0,\infty)\) 的总 Fisher 长度为
\[
\pi-2\arctan\sqrt{\varphi}.
\]
\end{theorem}
\begin{proof}
由定义
\[
\partial_q\log\pi_{m,q}(w)
=
\log d_m^{\mathrm{bin}}(w)-\partial_q\log S_q^{\mathrm{bin}}(m),
\]
故
\[
\mathcal I_m(q)
=
\operatorname{Var}_{\pi_{m,q}}\!\bigl(\log d_m^{\mathrm{bin}}(W)\bigr),
\qquad
W\sim \pi_{m,q}.
\]
另一方面，由定理 \ref{thm:fold-bin-two-state-asymptotic} 的一致对数展开，
\[
\log d_m^{\mathrm{bin}}(w)
=
m\log\!\Bigl(\frac{2}{\varphi}\Bigr)-w_m\log\varphi
+O\!\Bigl(\Bigl(\frac{\varphi}{2}\Bigr)^m\Bigr)
\]
在 \(w\in X_m\) 上一致成立。
再由定理 \ref{thm:xi-foldbin-escort-tv-collapse-onebit-gibbs}，
\[
\pi_{m,q}(w_m=1)
=
\theta(q)+O_q\!\Bigl(\Bigl(\frac{\varphi}{2}\Bigr)^m\Bigr).
\]
因此
\[
\mathcal I_m(q)
=
(\log\varphi)^2\operatorname{Var}_{\pi_{m,q}}(w_m)
+O_q\!\Bigl(\Bigl(\frac{\varphi}{2}\Bigr)^m\Bigr)
\]
并进一步得到
\[
\mathcal I_m(q)
=
(\log\varphi)^2\theta(q)(1-\theta(q))
+O_q\!\Bigl(\Bigl(\frac{\varphi}{2}\Bigr)^m\Bigr),
\]
从而
\[
\mathcal I_\infty(q)
=
(\log\varphi)^2\frac{\varphi^{q+1}}{(1+\varphi^{q+1})^2}.
\]

再计算
\[
u'(q)
=
2\cdot \frac{1}{1+\varphi^{q+1}}
\cdot \frac{\log\varphi}{2}\varphi^{(q+1)/2}
=
(\log\varphi)\frac{\varphi^{(q+1)/2}}{1+\varphi^{q+1}},
\]
故
\[
\bigl(u'(q)\bigr)^2
=
(\log\varphi)^2\frac{\varphi^{q+1}}{(1+\varphi^{q+1})^2}
=
\mathcal I_\infty(q).
\]
于是
\(
ds^2=\mathcal I_\infty(q)\,\dd q^2=du^2
\)，距离公式立得。
最后令 \(q_1=0\)、\(q_2\to\infty\)，便得到总长度
\(
\pi-2\arctan\sqrt{\varphi}
\)。
证毕。
\end{proof}
